\chapter{Namespaces, Types, and Include Paths}
\label{ch:sharp-runtime-namespaces}

Sharp-runtime makes a source-porting promise visible in the filesystem: if a developer knows
the .NET fully-qualified type name, the public include spelling should be predictable. The
module graph adds a second question that a header alone cannot answer: which component owns
that include and which dependencies follow transitively? This chapter establishes both rules,
then applies them to representative text, stream, encoding, and task types.

\section{Two public roots, no umbrella}

Exactly 1{,}032 public headers are reachable below two roots; five further \texttt{.hpp} files
are private implementation headers below module \texttt{src/} trees. \texttt{System/} contains
the literal .NET namespace tree; \texttt{SharpRuntime/} contains project-specific aliases and
source-porting helpers. There is no catch-all header. Private headers below a module's
\texttt{src/} directory are not added to consumer include paths.

The transliteration rule is mechanical:

\begin{lstlisting}
System.IO.FileStream
    -> #include "System/IO/FileStream.hpp"

System.Threading.Tasks.Task
    -> #include "System/Threading/Tasks/Task.hpp"

SharpRuntime::intcs
    -> #include "SharpRuntime/SharpRuntimeHelper.hpp"
\end{lstlisting}

One public type per like-named header is the norm. The rule makes a port review tractable: an
include that cannot be derived from the type name deserves inspection, and an accidental
include from another module's private tree cannot be hidden behind a universal umbrella.

\section{Namespaces do not equal components one-for-one}

The C++ namespace hierarchy describes API identity; the 44 CMake components describe build
ownership. Several namespace families split across components to keep link surfaces narrow.
Table~\ref{tab:sharp-namespace-components} gives representative mappings used by CNA.

\begin{table}[htbp]
  \centering
  \small
  \caption{Representative namespace/include/component mappings.}
  \label{tab:sharp-namespace-components}
  \begin{tabular}{@{}p{0.23\textwidth}p{0.43\textwidth}p{0.22\textwidth}@{}}
    \toprule
    API family & Include example & Direct component \\
    \midrule
    \texttt{System} core & \texttt{System/TimeSpan.hpp} & \texttt{Core.Base} \\
    Collections & \texttt{System/Collections/}\newline\texttt{Generic/List.hpp} & \texttt{Collections.Core} \\
    Object model & \texttt{System/Collections/}\newline\texttt{ObjectModel/Collection.hpp} & \texttt{Collections.}\newline\texttt{ObjectModel} \\
    Streams & \texttt{System/IO/MemoryStream.hpp} & \texttt{IO} \\
    Text & \texttt{System/Text/UTF8Encoding.hpp} & \texttt{Text} \\
    Tasks & \texttt{System/Threading/Tasks/Task.hpp} & \texttt{Threading.Tasks} \\
    Cryptographic hashes & \texttt{System/Security/}\newline\texttt{Cryptography/SHA256.hpp} & \texttt{Security.}\newline\texttt{Cryptography} \\
    \bottomrule
  \end{tabular}
\end{table}

The direct component is not the full closure. Its public dependencies carry lower-level
requirements, just as a C++ header's own includes do. This is why CNA names
\texttt{SharpRuntime::IO} rather than copying IO's current dependency list into its own CMake
files. Chapter~\ref{ch:sharp-components-consumption} covers the adapter and its legacy
monolith fallback.

\section{What CNA actually includes}

CNA directly names 54 distinct runtime headers. Fifty-two use the \texttt{System/} root. The
two \texttt{SharpRuntime/} spellings are \texttt{SharpRuntimeHelper.hpp}, which defines the
C\#-fidelity primitives and limits, and \texttt{Prop.hpp}, which exports property-declaration
macros for downstream code. Counting only the latter prefix and concluding that CNA consumes
two runtime headers mistakes a namespace prefix for a repository boundary.

\texttt{SharpRuntimeHelper.hpp} is unusually load-bearing: hundreds of public runtime headers
include it, and CNA uses aliases such as \texttt{intcs}, \texttt{bytecs}, and
\texttt{Single} throughout its public surface. It also exports two unnamespaced macros,
\texttt{CONTAINS} and \texttt{INTERNAL}; consumers must treat those as global preprocessor
names, not members protected by the \texttt{SharpRuntime} namespace.

\texttt{Prop.hpp} is the opposite kind of seam. The runtime's own property methods are
hand-written, and its modules do not include this macro header. CNA does. The file is therefore
an exported downstream facility even though a search limited to the library's implementation
would misleadingly classify it as unused.

\section{String and StringBuilder: a static-only utility versus a real mutable buffer}

\cnaclass{System::String} is deliberately \emph{not} a wrapper class at all --- its default
constructor and destructor are both explicitly deleted, making it a pure static-method utility
operating on ordinary \texttt{std::string} values (\texttt{IsNullOrEmpty}, \texttt{Compare},
the \texttt{IndexOf} / \texttt{LastIndexOf} family, \texttt{Split}, \texttt{Substring},
\texttt{Trim*}, \texttt{PadLeft} / \texttt{PadRight}). \cnaclass{Text::StringBuilder}, by
contrast, is a real, stateful, mutable buffer type --- the many-overloaded \texttt{Append}
(string, char, int, double, float, bool, and the \texttt{longcs} alias from
Chapter~\ref{ch:design-philosophy}), \texttt{AppendLine}, \texttt{Insert},
\texttt{Remove}, \texttt{Replace}, and a nested \texttt{ChunkEnumerator} class mirroring real
.NET's own chunked-buffer enumeration shape --- real .NET's own \texttt{StringBuilder} is
internally a linked list of separate character chunks, and \texttt{ChunkEnumerator} walks that
list without copying; this port's \texttt{StringBuilder} is backed by a single
\texttt{std::string} instead, so its own \texttt{ChunkEnumerator} always yields exactly one
chunk holding the entire current buffer --- the enumeration \emph{shape} is preserved
faithfully even though the simpler single-buffer implementation underneath never actually
produces more than one chunk to walk.

\begin{lstlisting}
// Worked example: System::String is a pure static utility -- there is no
// instance to construct, only free functions operating on std::string.
std::vector<std::string> parts = System::String::Split(
    "Content/Textures/hero.png", '/');
std::string trimmed = System::String::Trim(parts.back());

// Text::StringBuilder, by contrast, is a real mutable buffer.
System::Text::StringBuilder sb;
sb.Append("Score: ").Append(playerScore).AppendLine();
sb.Append("Lives: ").Append(livesRemaining);

for (const std::string& chunk : sb.GetChunks())
{
    // Exactly one iteration, always -- see the ChunkEnumerator note above.
    RenderHudText(chunk);
}
\end{lstlisting}

\section{Stream: a lightweight subset, honest about it}

\cnaclass{IO::Stream}'s own class comment states its scope plainly: ``lightweight subset of
the .NET Stream API.'' \texttt{Read(buffer, offset, count)}, \texttt{Close()}, and
\texttt{getLengthProperty()} are pure virtual; \texttt{Write} / \texttt{WriteByte} default to
throwing \texttt{NotSupportedException} unless a writable stream subclass overrides them ---
the same shape real .NET's own read-only stream implementations follow, rather than forcing
every stream subclass to implement a full read/write/seek contract regardless of whether it
actually supports all three. \texttt{getCanWriteProperty()} defaults to false and
\texttt{getCanReadProperty()} defaults to true, matching a read-only stream's honest self-report
without requiring every subclass to override both just to state the obvious.

\begin{lstlisting}
// Worked example: the minimal real subclass this base class actually
// requires -- only three overrides, matching a genuinely read-only
// stream's real capabilities (Write/WriteByte are correctly inherited as
// NotSupportedException-throwing, never called by a well-behaved reader).
class MemoryBlobStream final : public System::IO::Stream
{
public:
    explicit MemoryBlobStream(std::vector<SharpRuntime::bytecs> data)
        : data_(std::move(data)) {}

    SharpRuntime::intcs Read(SharpRuntime::bytecs buffer[],
                             SharpRuntime::intcs offset, SharpRuntime::intcs count) override
    {
        if (buffer == nullptr && count != 0)
            throw System::ArgumentNullException("buffer");
        if (offset < 0)
            throw System::ArgumentOutOfRangeException("offset");
        if (count < 0)
            throw System::ArgumentOutOfRangeException("count");

        SharpRuntime::intcs n = std::min<SharpRuntime::intcs>(
            count, static_cast<SharpRuntime::intcs>(data_.size()) - position_);
        if (n != 0)
            std::memcpy(buffer + offset, data_.data() + position_, n);
        position_ += n;
        return n;
    }

    void Close() override {}

    [[nodiscard]] SharpRuntime::intcs getLengthProperty() const override
    {
        return static_cast<SharpRuntime::intcs>(data_.size());
    }

private:
    std::vector<SharpRuntime::bytecs> data_;
    SharpRuntime::intcs position_ = 0;
};
\end{lstlisting}

\subsection{\cnaclass{MemoryStream}: writable by default, read-only when requested}

\cnaclass{System::IO::MemoryStream} is the concrete \texttt{Stream} subclass this ecosystem
actually uses in practice --- CNA's own \texttt{ContentManager} reaches for it to wrap data
loaded from sources like Android APK assets, where a real filesystem path is not available at
all. Its default constructor creates an empty, writable buffer, matching real .NET. The
buffer-copy constructor now makes the same default explicit while retaining a C\texttt{++}-only
way to request a read-only copy:

\begin{lstlisting}
MemoryStream(const bytecs* buffer, intcs size, bool writable = true);
\end{lstlisting}

The implementation validates before forming the vector range. A negative size is reported as
\cnaclass{ArgumentOutOfRangeException}; a null pointer with non-zero size is reported as
\cnaclass{ArgumentNullException}. The one deliberate source-language adaptation is
\texttt{nullptr, 0}, accepted as an empty range because C\texttt{++} can express an empty pointer
range whereas .NET's array overload simply rejects a null array. Passing \texttt{false} creates
a readable, seekable, non-writable copy whose \texttt{Write} and \texttt{WriteByte} calls throw
\cnaclass{NotSupportedException}; omitting it matches the writable .NET array constructor.

\subsection{Closing a memory stream closes operations, not the bytes it already owns}

The constructor policy above should not be confused with this class's close/dispose behavior.
\texttt{MemoryStream::Close()} changes only its
\texttt{isOpen\_} state; it deliberately retains both the vector of bytes and the current
position.  After close, \texttt{Read}, \texttt{Write}, \texttt{WriteByte}, \texttt{Seek},
\texttt{Length}, and \texttt{Position} correctly reject use with
\cnaclass{ObjectDisposedException}, while \texttt{CanRead} and \texttt{CanSeek} return false.
A closed stream created as writable deliberately retains \texttt{CanWrite == true}: .NET's own
property reports the construction-time writable flag even though subsequent writes throw.
In further deliberate contrast, \texttt{ToArray()} and \texttt{GetBuffer()} remain usable:

\begin{lstlisting}
System::IO::MemoryStream stream;
const SharpRuntime::bytecs payload[] = {1, 2, 3, 4};
stream.Write(payload, 0, 4);
stream.Close();

std::vector<SharpRuntime::bytecs> saved = stream.ToArray();  // still {1,2,3,4}
// stream.getLengthProperty();  // throws ObjectDisposedException
// stream.Read(buffer, 0, 1);   // throws ObjectDisposedException
\end{lstlisting}

That asymmetry is intentional, not a use-after-dispose loophole. Real .NET explicitly preserves
the buffer in its own dispose implementation, \texttt{MemoryStream::Dispose(bool)}, so
post-disposal \texttt{TryGetBuffer}, \texttt{GetBuffer}, and \texttt{ToArray} can still retrieve data; the
runtime's \texttt{ensureNotClosed()} guard follows the same rule by guarding operations but not
the two extraction methods.  CNA callers can therefore finish producing an in-memory content
blob, close the stream to catch accidental later I/O, and still hand a stable copy to a consumer.
The distinction between the methods matters: \texttt{ToArray()} returns an independent vector,
safe across later stream writes, whereas \texttt{GetBuffer()} is a reference to the live vector
and may be invalidated by a reallocation.

The same implementation also has two less obvious safety properties whose tests prevent a
``memory only'' stream from becoming an unchecked byte container. Null buffers produce
\cnaclass{ArgumentNullException}; negative offsets or counts produce
\cnaclass{ArgumentOutOfRangeException}. Neither case is mistaken for an EOF return of zero.
And a caller may legally set \texttt{Position} past the current end, but a later write computes
\texttt{position + count} in \texttt{int64\_t} before resizing.  If that sum exceeds the
runtime's signed \texttt{intcs} range it throws \cnaclass{IO::IOException} (``Stream was too
long.'') instead of wrapping negative and writing beyond the vector.  This was not defensive
theorizing: the former narrow-integer calculation had the same signed-overflow shape already
reproduced under UBSan in a related slicing path, and \texttt{StreamTests.cpp} now exercises
\texttt{Position = 2147483647} followed by a ten-byte write specifically to require the
exception rather than an allocation or memory corruption.

\section{\cnaclass{UTF8Encoding}: conformance validation and fallback evidence}

\cnaclass{Text::UTF8Encoding}'s own class comment states a real, specific guarantee that is
easy to assume any UTF-8 encoder provides and that this port genuinely earns rather than
skips: \texttt{GetBytes()} / \texttt{GetString()} \emph{validate} that input is well-formed
UTF-8 --- not merely pass bytes through as opaque \texttt{char} values --- and route any
ill-formed byte at each invalid position through the configured fallback before resuming.
Reading \texttt{UTF8Encoding.cpp}'s own local \texttt{wellFormedUtf8Length()} helper directly
shows this is real conformance checking, not a length-only byte count: a leading byte is
walked against the real UTF-8 grammar (1/2/3/4-byte forms, each continuation byte checked for
the \texttt{10xxxxxx} pattern), and even a \emph{structurally} well-formed sequence is rejected
if the code point it decodes to is an overlong encoding (a multi-byte sequence encoding a code
point that a shorter form could already represent --- a real, historically security-relevant
UTF-8 attack shape, since a naive decoder that accepts overlong encodings lets an attacker
smuggle a byte value like \texttt{0x2F} (\texttt{'/'}) past a filter checking only the raw byte
stream) or a UTF-16 surrogate code point (\texttt{U+D800}--\texttt{U+DFFF}, which UTF-8 must
never directly encode, since surrogates exist only as a UTF-16 encoding artifact).

\begin{lstlisting}
// From wellFormedUtf8Length() -- the 3-byte case, real conformance logic,
// not just "are there two continuation bytes after this lead byte":
if ((c0 & 0xF0) == 0xE0 && /* two real continuation bytes follow */) {
    uint32_t cp = /* decode the 16-bit code point from all three bytes */;
    return (cp >= 0x800 && !(cp >= 0xD800 && cp <= 0xDFFF)) ? 3 : 0;
    //      ^ rejects overlong (cp < 0x800 has a shorter valid encoding)
    //                          ^ rejects UTF-16 surrogates directly encoded in UTF-8
}
\end{lstlisting}

The constructor also earns a real, specific correction over the naive default: real .NET's own
\texttt{UTF8Encoding.SetDefaultFallbacks()} substitutes the Unicode replacement character
\texttt{U+FFFD} for invalid input, not the generic \cnaclass{Encoding} base class's plain
\texttt{"?"} default (correct only for single-byte code pages, where every valid character
already has an unambiguous byte representation to fall back to). This port's constructor
reproduces that distinction precisely:

\begin{lstlisting}
UTF8Encoding::UTF8Encoding() {
    // Real UTF8Encoding.SetDefaultFallbacks() uses a U+FFFD replacement
    // fallback, not the generic Encoding base class's "?" default.
    setEncoderFallbackProperty(
        std::make_shared<EncoderReplacementFallback>("\xEF\xBF\xBD"));
    setDecoderFallbackProperty(
        std::make_shared<DecoderReplacementFallback>("\xEF\xBF\xBD"));
}
\end{lstlisting}

A worked example makes the whole mechanism concrete --- three real byte sequences, one
well-formed and two deliberately malformed in different ways, through the same
\texttt{GetString()} call:

\begin{lstlisting}
System::Text::UTF8Encoding utf8;

// Well-formed: a real 3-byte sequence for U+20AC (EURO SIGN).
const SharpRuntime::bytecs euroSign[] = {0xE2, 0x82, 0xAC};
std::string ok = utf8.GetString(euroSign, 0, 3);          // EURO SIGN, unchanged

// Overlong encoding of U+002F ('/') using the 2-byte form -- structurally
// valid continuation-byte shape, but wellFormedUtf8Length() rejects it
// because cp (0x2F) is below the 2-byte form's own 0x80 minimum.
const SharpRuntime::bytecs overlongSlash[] = {0xC0, 0xAF};
std::string rejected = utf8.GetString(overlongSlash, 0, 2);  // two U+FFFD

// A lone continuation byte with no lead byte before it -- fails every one
// of wellFormedUtf8Length()'s four branches, falls through to the final
// "return 0" unconditionally.
const SharpRuntime::bytecs loneContinuation[] = {0x80};
std::string alsoRejected = utf8.GetString(loneContinuation, 0, 1);  // one U+FFFD
\end{lstlisting}

The overlong and lone-continuation cases each substitute one \texttt{U+FFFD} per invalid
\emph{byte}, not per attempted sequence --- \texttt{overlongSlash}'s two bytes both fail
validation independently (the loop advances one byte at a time through the fallback path, per
\texttt{UTF8Encoding.cpp}'s own \texttt{++i} in the non-well-formed branch), so a 2-byte
overlong sequence becomes two replacement characters, not one. This matches real .NET's own
observable behavior for the identical input, confirmed by working through the .NET runtime's
own documented per-byte fallback-and-resynchronize algorithm for UTF-8 decoding rather than
assumed by analogy.

This behavior now has direct regression evidence. \texttt{UTF8EncodingTests} covers malformed
continuations, overlong input, truncated sequences, directly encoded surrogates and exception
fallbacks, including the one-replacement-per-byte resynchronization used above. The separate
\texttt{Utf8Tests} family validates overlong sequences, surrogates, code points above
\texttt{U+10FFFF}, truncation and lone continuation bytes through the span-oriented validator.
The source read explains the algorithm; these two distinct test surfaces pin both the encoding
facade's observable substitutions and the lower-level validity classification.

\section{Generic continuations exist; generic combinators do not}

Both non-generic \cnaclass{Task} and generic-result \texttt{TaskT<TResult>} have real
\texttt{ContinueWith} implementations. Their shared-state shapes each contain a mutex,
condition variable, and callback list; registered callbacks execute inline on the thread that
completes the antecedent, or synchronously on the caller when it is already complete.
\texttt{TaskT<TResult>} provides both an action overload returning \cnaclass{Task} and a
result-producing overload returning \texttt{TaskT<TNewResult>}. The tests cover success, fault,
cancellation inspection, filtered options, result chaining, empty callables, and release of
captured state. Both implementations capture the antecedent state weakly to avoid a cycle from
that state's own continuation list back to itself.

The narrower missing surface is aggregation over generic tasks. Static \texttt{WhenAll} and
\texttt{WhenAny} accept only \texttt{std::vector<Task>}; there is no corresponding overload for
\texttt{std::vector<TaskT<TResult>>} that produces a result array or a winning generic task.
Thus callers can attach a continuation to one generic result, but cannot express .NET's generic
multi-task combinators through this API. A second minor boundary is that the lightweight
antecedent reconstructed for a callback shares terminal status/result/exception state but not
the original separately-held cancellation-token object, so its token property reports
\texttt{None}.

\section{Reading a dependency from source}

A reliable dependency review uses three layers in order. First derive the public include from
the type's fully-qualified name. Then locate the module whose public include directory owns the
file. Finally inspect that component's declared public dependencies rather than inferring a
closure from nested namespaces. Namespace depth, directory depth, and link depth often align,
but the component registry --- particularly compatibility umbrellas and the \texttt{Xml.XPath}
alias --- is the authority when they do not.

The representative types in this chapter also show why include availability is weaker than
semantic parity. \cnaclass{UTF8Encoding}'s malformed-input behavior is implemented and directly
tested; \cnaclass{TaskT<TResult>} has generic continuations but not generic multi-task
combinators. A mechanically correct path proves that a type can be named. It does not prove that
every .NET behavior behind that name exists.
